\documentclass[a4paper,10pt]{article}
\usepackage[margin=1.2cm]{geometry}
\usepackage[T1]{fontenc}
\usepackage[utf8]{inputenc}
\usepackage{lmodern}
\usepackage[hidelinks]{hyperref}
\pagestyle{empty}
\setlength{\parindent}{0pt}
\setlength{\parskip}{2pt}
\begin{document}
\textbf{Israel Souza Ferreira}\par
\textbf{Data Scientist}\par
Fortaleza, CE, Brasil · Remoto\par
israel.souza.dev@gmail.com | +55 85 98433-2790\par
https://www.linkedin.com/in/isrreal/ | https://github.com/isrreal\par
https://israelsouza.speculummaius.com.br/\par
\vspace{5pt}\textbf{Resumo Profissional}\par
Bacharel em Ciência da Computação com experiência em P\&D de Machine Learning multimodal, representação de dados e avaliação de modelos. Combino modelagem, experimentação quantitativa e implementação de sistemas para investigar hipóteses.\par
\vspace{5pt}\textbf{Experiência Profissional}\par
\textbf{Tieta.ai · P\&D em Machine Learning · Remoto}\par
Pesquisador de P\&D: 07/2026–presente | Pesquisador de ML, CNPq RHAE (meio período): 04/2025–06/2026\par
P\&D colaborativo em Machine Learning multimodal e representation learning para análise educacional, com experimentação de modelos e componentes baseados em LLMs.\par
\hangindent=1em\textbullet\ Colaborei na investigação de hipóteses, implementação de baselines e análise quantitativa de experimentos, com apoio da literatura científica para orientar a modelagem.\par
\hangindent=1em\textbullet\ Participei da implementação e avaliação de modelos de deep learning multimodal, comparando alternativas de representação e combinação de modalidades.\par
\hangindent=1em\textbullet\ Colaborei no desenvolvimento e na avaliação de representações vetoriais e recuperação de informação, investigando embeddings e medidas de similaridade para texto e imagem.\par
\hangindent=1em\textbullet\ Atuei na experimentação de componentes baseados em modelos de linguagem (LLMs), com foco em extração estruturada de informação e representação semântica.\par
\textbf{Consultoria autônoma · Instituição de ensino · 03/2026–presente}\par
Desenvolvedor Full Stack e Engenheiro de ML · Face Clock Evoluir\par
v1.3.0-rc.3: candidata a piloto institucional acompanhado, implantada em ambiente de testes.\par
\hangindent=1em\textbullet\ Comparei baselines clássicos, neurais e híbridos para triagem documental: F1 0,778 e acurácia 0,875 no teste reservado de 64 documentos; amostra pequena limita generalização.\par
\hangindent=1em\textbullet\ Recalibrei o limiar facial: desconhecidos aceitos passaram de 45,4\% a 0,2\% em teste controlado. Calibração com a população da instituição permanece pendente.\par
\vspace{5pt}\textbf{Pesquisa e Projetos}\par
\textbf{Dominação Romana Tripla em Grafos · TCC · 2024–2025}\par
\hangindent=1em\textbullet\ Identifiquei contraexemplo de uma formulação PLI publicada e implementei formulação corrigida, GA e ACO (C++, Python/Pyomo). Avaliei 362 grafos; uma execução por instância limita as conclusões.\par
\textbf{Auxílio Emergencial · Em desenvolvimento}\par
\hangindent=1em\textbullet\ Implementei ingestão em blocos com Pandas e PostgreSQL COPY. Profiling em recortes de até 5 milhões de linhas mostrou memória crescente; a correção por staging ainda não foi validada.\par
\vspace{5pt}\textbf{Competências Técnicas}\par
Validação cruzada, calibração, ablação, métricas de classificação; C++, GA, ACO, Python/Pyomo.\par
Python, PyTorch, scikit-learn; deep learning multimodal, representation learning, embeddings, recuperação de informação, LLMs.\par
FastAPI, PostgreSQL/pgvector, Docker, ONNX Runtime, Git, pytest; SQL, Pandas, NumPy.\par
\vspace{5pt}\textbf{Formação Acadêmica}\par
\textbf{Bacharelado em Ciência da Computação · Universidade Federal do Ceará (UFC), Quixadá}\par
2021–19/08/2026 · TCC defendido em 02/2025.\par
\vspace{5pt}\textbf{Idiomas}\par
Português nativo. Inglês: boa comunicação; leitura técnica fluente de artigos e documentação.\par
\end{document}
